\subsection{Inverse Function Theorem}

% **Role:** calculus as local linearization plus global geometric integration.
%
% Inverse function theorem, implicit function theorem, rank theorem, local normal forms.

\begin{theorem} \label{thm:finite-dimensional-linear-maps-lipschitz-homeomorphisms} \hlt{Finite-Dimensional Linear Maps Are Lipschitz}\\
Let $A\in\LL(\R^n,\R^m)$.
\begin{enumerate}[label=(\roman*)]
\setlength{\itemsep}{0pt}
\item For all $\x,\y\in\R^n$, $\norm{A\x-A\y}_2\leq\normop{A}\norm{\x-\y}_2$; hence $A$ is $\normop{A}$-Lipschitz.
\item If $m=n$ and $A$ is invertible, then $\frac{1}{\normop{A^{-1}}}\norm{\x-\y}_2\leq\norm{A\x-A\y}_2\leq\normop{A}\norm{\x-\y}_2$ for all $\x,\y\in\R^n$; hence $A:\R^n\to\R^n$ is a homeomorphism.
\end{enumerate}
\begin{proof}
\begin{enumerate}[label=(\roman*)]
\setlength{\itemsep}{0pt}
\item Definition \ref{def:operator-norm-rn-rm} gives $\norm{A\z}_2\leq\normop{A}\norm{\z}_2$ for all $\z\in\R^n$; set $\z=\x-\y$.
\item Apply (i) to $A^{-1}$: $\norm{\x-\y}_2=\norm{A^{-1}(A\x-A\y)}_2\leq\normop{A^{-1}}\norm{A\x-A\y}_2$.\\
This gives the lower bound; the upper bound is (i). Hence $A$ and $A^{-1}$ are continuous.
\end{enumerate}
\hlt{Structural Commentary}: $A^{-1}$ converts upper control into lower control.\\
\hlt{Key Insight}: Invertible linear maps are bi-Lipschitz.
\end{proof}
\end{theorem}

\begin{theorem} \label{thm:invertible-operators-open-and-inversion-continuous} \hlt{Invertible Operators Form an Open Set}\\
Let $\Omega$ be the set of invertible operators in $\LL(\R^n)$.
\begin{enumerate}[label=(\roman*)]
\setlength{\itemsep}{0pt}
\item If $A\in\Omega$, $B\in\LL(\R^n)$, and $\normop{B-A}\,\normop{A^{-1}}<1$, then $B\in\Omega$.
\item $\Omega$ is open in $\LL(\R^n)$.
\item The inversion map $\iota:\Omega\to\Omega$, $\iota(A)=A^{-1}$, is continuous.
\end{enumerate}
\begin{proof}
\begin{enumerate}[label=(\roman*)]
\setlength{\itemsep}{0pt}
\item Put $\alpha\defeq1/\normop{A^{-1}}$ and $\beta\defeq\normop{B-A}<\alpha$.\\
For every $\x\in\R^n$, $\alpha\norm{\x}_2\leq\norm{A\x}_2\leq\beta\norm{\x}_2+\norm{B\x}_2$, hence $(\alpha-\beta)\norm{\x}_2\leq\norm{B\x}_2$.\\
Thus $\ker B=\{\ze\}$, so $B\in\Omega$ by Theorem \ref{thm:finite-dimensional-operator-invertible-equivalences}; setting $\x=B^{-1}\y$ gives $\normop{B^{-1}}\leq(\alpha-\beta)^{-1}$.
\item For $A\in\Omega$, the ball $\{B:\normop{B-A}<1/\normop{A^{-1}}\}$ lies in $\Omega$ by (i); hence $\Omega$ is open.
\item If $B\to A\in\Omega$, then eventually $\beta<\alpha/2$, so $B\in\Omega$ and $\normop{B^{-1}}\leq(\alpha-\beta)^{-1}\leq2/\alpha$.\\
Since $B^{-1}-A^{-1}=B^{-1}(A-B)A^{-1}$, $\normop{B^{-1}-A^{-1}}\leq(2/\alpha)\normop{B-A}\normop{A^{-1}}\to0$.
\end{enumerate}
\hlt{Structural Commentary}: Perturbations below the lower stretch stay invertible.\\
\hlt{Key Insight}: $B^{-1}-A^{-1}$ factors through $A-B$.
\end{proof}
\end{theorem}

\begin{perspective} \label{per:linear-inverse-lower-bound} \hlt{Linear Inverse Estimate}\\
For invertible $A\in\LL(\R^n)$, Theorem \ref{thm:finite-dimensional-linear-maps-lipschitz-homeomorphisms} gives $c\norm{\h}_2\leq\norm{A\h}_2\leq\normop{A}\norm{\h}_2$; and $c\defeq\frac{1}{\normop{A^{-1}}}$.\\
The inverse function theorem preserves this lower stretch under small derivative changes.
\end{perspective}

\begin{figure}[H]
  \centering
  \begin{tikzpicture}[
  scale=0.95,
  >=Stealth,
  every node/.style={font=\scriptsize},
  ambient/.style={draw=black!35, fill=gray!5, line width=0.55pt},
  image/.style={draw=blue!70!black, fill=blue!5, line width=0.8pt},
  point/.style={circle, fill=black, inner sep=1.2pt},
  distline/.style={<->, blue!70!black, line width=0.75pt}
]
  \draw[ambient] (-2.20,-1.25) rectangle (3.00,1.55);
  \node[anchor=north west] at (-2.00,1.36) {$W$};

  \coordinate (O) at (-1.45,0);
  \coordinate (Q) at (-0.25,0);
  \coordinate (C) at (1.00,0);

  \draw[image] (C) ellipse (1.25 and 0.58);

  \node[point] at (O) {};
  \node[point] at (Q) {};

  \node[below=3pt] at (O) {$0$};
  \draw[distline] (O) -- node[above=2pt] {$c$} (Q);
  \node[above=4pt] at (1.00,0.58) {$A(S^{n-1})$};
\end{tikzpicture}

  \caption{Invertibility keeps $A(S^{n-1})$ a positive distance from $0$.}
  \label{fig:linear-isomorphism-conorm}
\end{figure}

\begin{theorem} \label{thm:inverse-function-theorem-rn} \hlt{Inverse Function Theorem}\\
Let $E\subseteq\R^n$ be open, let $\f:E\to\R^n$ be $C^1$, let $\a\in E$, put $\b\defeq\f(\a)$, and assume $A\defeq D\f(\a)$ is invertible. Then there exist open neighborhoods $U$ of $\a$ and $V$ of $\b$ such that:
\begin{enumerate}[label=(\roman*)]
\setlength{\itemsep}{0pt}
\item $\f|_U:U\to V$ is bijective; write $\g\defeq(\f|_U)^{-1}:V\to U$.
\item $\g$ is differentiable on $V$ and $D\g(\y)=\bigl(D\f(\g(\y))\bigr)^{-1}$ for all $\y\in V$.
\item $\g\in C^1(V,\R^n)$; equivalently, $\f|_U$ is a $C^1$ diffeomorphism.
\end{enumerate}
\begin{proof} Choose $0<\theta<1$ and $r>0$ so that $\overline{B}_r(\a)\subset E$ and $\normop{I-A^{-1}D\f(\x)}\leq\theta$ on $\overline{B}_r(\a)$.\\
Let $V\defeq\{\y:\norm{A^{-1}(\y-\b)}_2<(1-\theta)r\}$.
\begin{enumerate}[label=(\roman*)]
\setlength{\itemsep}{0pt}
\item For $\y\in V$, set $\varphi_{\y}(\x)\defeq\x+A^{-1}(\y-\f(\x))$. Theorem \ref{thm:bounded-derivative-convex-mean-value} gives $\norm{\varphi_{\y}(\x)-\varphi_{\y}(\mathbf{z})}_2\leq\theta\norm{\x-\mathbf{z}}_2$ on $\overline{B}_r(\a)$, and $\norm{\varphi_{\y}(\x)-\a}_2\leq\theta r+\norm{A^{-1}(\y-\b)}_2<r$. Since $\overline{B}_r(\a)$ is complete, Theorem \ref{thm:contraction-principle} gives a unique fixed point $\x\in\overline{B}_r(\a)$, equivalently $\f(\x)=\y$. With $U\defeq B_r(\a)\cap\f^{-1}(V)$, $\f|_U:U\to V$ is bijective.
\item Let $\y=\f(\x)$ and $\y+\mathbf{k}=\f(\x+\h)$ lie in $V$. The derivative bound gives $\norm{\h-A^{-1}\mathbf{k}}_2\leq\theta\norm{\h}_2$, hence $\norm{\h}_2\leq C\norm{\mathbf{k}}_2$ locally. Since $\normop{I-A^{-1}D\f(\x)}<1$, $D\f(\x)$ is invertible, and\\
$\g(\y+\mathbf{k})-\g(\y)-\bigl(D\f(\x)\bigr)^{-1}\mathbf{k}=-\bigl(D\f(\x)\bigr)^{-1}[\f(\x+\h)-\f(\x)-D\f(\x)\h]$.\\
Since $\f$ is differentiable at $\x$, division by $\norm{\mathbf{k}}_2$ gives $D\g(\y)=\bigl(D\f(\x)\bigr)^{-1}$.
\item The estimate in (ii) makes $\g$ continuous. Since $D\g(\y)=\bigl(D\f(\g(\y))\bigr)^{-1}$, continuity of $D\g$ follows from continuity of $D\f$, $\g$, and Theorem \ref{thm:invertible-operators-open-and-inversion-continuous}.
\end{enumerate}
\hlt{Structural Commentary}: $A^{-1}$ turns solving $\f(\x)=\y$ into a contraction.\\
\hlt{Key Insight}: Small derivative variation preserves invertibility locally.
\end{proof}
\end{theorem}

\begin{corollary} \label{cor:local-diffeomorphism-open-map} \hlt{Local Diffeomorphisms Are Open}\\
If $\f:E\subseteq\R^n\to\R^n$ is $C^1$ and $D\f$ is everywhere invertible, then $\f$ is an open map (locally one-to-one).
\begin{proof}
Let $W\subseteq E$ be open and let $\y\in\f(W)$. Choose $\x\in W$ with $\f(\x)=\y$. Applying Theorem \ref{thm:inverse-function-theorem-rn} to $\f|_W$ at $\x$ gives open neighborhoods $U\subseteq W$ of $\x$ and $V$ of $\y$ with $\f(U)=V$. Hence $V\subseteq\f(W)$, so $\y$ is interior to $\f(W)$. Since $\y$ was arbitrary, $\f(W)$ is open. Local one-to-one-ness is the same theorem applied pointwise.\\
\hlt{Structural Commentary}: Openness is checked one image point at a time.\\
\hlt{Key Insight}: A local inverse turns nearby domain points into a full image neighborhood.
\end{proof}
\end{corollary}
